\section{CSP、删域与局部修复}\label{sec:csp}
\src{课件 06 p.14--63；HW4 CSP1--3}
CSP 由变量、域、约束组成；解为\textbf{完整且满足全部约束}的赋值。
\par\noindent\begin{tabularx}{\textwidth}{lY}
\toprule 名称&执行规则\\\midrule
MRV&选当前合法域最小的未赋值变量（不是最小数字值）。\\
Degree&通常用来打破 MRV 并列，选与最多未赋值变量有约束者。\\
LCV&对已选变量，选删掉其他域中值最少的候选值。\\
Forward checking&赋值后删未赋值邻居中不相容值，任一空域则回溯。\\
Arc consistency&每个源域值在目标域有支持；没支持则删源域值，继续传播。\\
Min-conflicts&从完整但有冲突赋值出发，选冲突变量，把它改成违反最少约束的值。\\\bottomrule
\end{tabularx}\par
回溯返回上一层时要恢复该层造成的删域。

\ex{建模公式：皇后与密码算术}
皇后每列一个：$Q_i\in\{1,\ldots,n\}$；对 $i<j$，$Q_i\ne Q_j$，$|Q_i-Q_j|\ne|i-j|$。候选数 $n^n$；若预设行也互异才是 $n!$。解例 4 皇后 $(2,4,1,3)$。

TWO+TWO=FOUR：六字母 AllDifferent，数字 0--9，首位 T,F 非零，进位 $X_1,X_2,X_3\in\{0,1\}$：
\[
2O=R+10X_1,\quad X_1+2W=U+10X_2,\quad
X_2+2T=O+10X_3,\quad X_3=F=1.
\]
进位不属于字母 AllDifferent，不额外添加不同进位必须不同的条件。

\ex{MRV + degree + FC 全表（HW4 CSP1）}
边 $E=\{ab,ac,ad,bc,be,cd,ce,de\}$，端点不等。先做一元约束：$D_a=\{1,2\},D_b=\{1,2,3\}$，其余 $\{1,2,3,4\}$。规则 MRV→未赋值 degree→字母序；值升序。
\begin{center}\small\begin{tabular}{c|c|ccccc}\toprule
步&赋值&$D_a$&$D_b$&$D_c$&$D_d$&$D_e$\\\midrule
0&--&12&123&1234&1234&1234\\
1&a=1&--&23&234&234&1234\\
2&b=2&--&--&34&234&134\\
3&c=3&--&--&--&24&14\\
4&d=2&--&--&--&--&14\\
5&e=1&--&--&--&--&--\\\bottomrule\end{tabular}\end{center}
表中 24 表示集合 $\{2,4\}$；横线表示已赋值。答案 $\ans{(a,b,c,d,e)=(1,2,3,2,1)}$，无需回溯。d=4 也能完成，但不是本例数值升序选择。
\checkit{逐条检查八条边，不存在 b--d 边，所以 b=d=2 允许。a 与 e 不相邻，均为 1 允许。}

\ex{LCV、弧一致性与 AC-3（教学变式）}\label{ex:ac}
\textbf{LCV：}X 域 12，与 Y 域 1、Z 域 12 都要求不同。X=1 删 Y 的 1（空域）和 Z 的 1；X=2 只删 Z 的 2，故选 2。

\textbf{弧方向：}X$\to$Y 检查 $\forall x\in D_X,\exists y\in D_Y:C(x,y)$。若 X$\ne$Y，$D_X=\{R,B\},D_Y=\{B\}$，X 的 B 无支持，删 B 得 $D_X=\{R\}$。

\textbf{传播：}A$\ne$B、B$\ne$C，域 $D_A=\{R\},D_B=\{R,B\},D_C=\{B\}$。B$\to$A 删红，B 只剩蓝；C$\to$B 删蓝，C 空，当前分支无解。

AC-3：队列存有向弧；取 X$\to$Y 做 REVISE；若 $D_X$ 变化，把其他邻居的 Z$\to$X 重新排队；直到空域失败或队列空。常见复杂度 $O(ed^3)$。
\checkit{只在支持不存在时删值；弧一致不等于全局可解。两色三角形每条弧都有支持，整体仍无解。}

\ex{固定红点染色与四色下界（HW4 CSP2--3）}
CSP2：左 L，上左 A、下左 B、上右 C、下右 D；边 LA,LB,AB,AC,AD,BD,CD。L=红，故 A、B 必为绿蓝且不同；取 A=绿，B=蓝，D 被 A、B 排除只能红，C 被 A、D 排除只能蓝。
\[
\ans{(L,A,B,C,D)=(R,G,B,B,R)}.
\]
CSP3：修改地图的 SA,Q,NSW,V 形成 $K_4$，至少四色；一组可行值 WA=2,NT=3,SA=1,Q=2,NSW=3,V=4，所以 $\chi=4$。下界需要 K4 或等价排除理由；只画四色方案只能证明至多四。

\ex{min-conflicts 冲突计数（课件 06 p.54--58；教学变式）}
4 皇后每列一枚，初始 $(1,1,1,1)$。选第 2 列，候选行 1,2,3,4 与其他三枚皇后的冲突数分别 $3,2,1,0$，移到行 4 得 $(1,4,1,1)$。这是一次修复，尚非解；继续直到每对满足行/对角约束。整盘冲突按\textbf{皇后对}计一次，不能把各皇后冲突数相加后不除 2。

\ex{树结构与 Zebra 表示简查}
树状二元 CSP：叶到根做父对子的弧一致，再根到叶选择与父相容值；非空域时可解，典型 $O(nd^2)$。含环图可 cutset 固定后解树，或簇分解后在共享变量上匹配；簇大时组合数仍指数。

Zebra 位置变量取 1--5，同类 AllDifferent；同房用等号，邻房用 $|pos(A)-pos(B)|=1$，绿紧接象牙右侧为 $pos(G)=pos(I)+1$。挪威人在1、邻蓝故蓝在2；第一间排除红/蓝/绿/象牙后黄；黄配 Kit Kats，再排除牛奶/咖啡/茶/橙汁，故挪威人喝水。
